\documentclass[11pt]{amsart}
\usepackage[margin=1in]{geometry}
\usepackage{amsmath,amssymb,amsthm,mathtools,booktabs}
\usepackage{enumitem}
\usepackage{hyperref}
\hypersetup{colorlinks=true,linkcolor=blue,citecolor=blue,urlcolor=blue}
\usepackage[final]{microtype}
\setlength{\emergencystretch}{3em}

\newcommand{\leanrepo}{https://github.com/HQIV/hqiv-lean/blob/main}
\newcommand{\leanfile}[1]{\href{\leanrepo/#1}{\nolinkurl{#1}}}
\newcommand{\leanid}[1]{\nolinkurl{#1}}

\title[SM Lagrangian from HQIV Discrete Action]{Building the Standard Model Lagrangian from\\
the HQIV Discrete Action: A Sector-by-Sector Lean Bridge}

\author{Steven Ettinger Jr}
\address{Independent Researcher}
\email{steven@disregardfiat.tech}
\date{Preprint draft --- \today}

\newtheorem{theorem}{Theorem}[section]
\newtheorem{lemma}[theorem]{Lemma}
\newtheorem{proposition}[theorem]{Proposition}
\newtheorem{definition}[theorem]{Definition}
\newtheorem{remark}[theorem]{Remark}
\newtheorem{corollary}[theorem]{Corollary}

\begin{document}

\begin{abstract}
\noindent\textbf{Position in the publication chain.}
This is the fifth paper in the HQIV preprint series and the
second of the Tier-2 Standard-Model layer. Tier~0 fixes the
framework \cite{ettinger2026hqiv,hqiv-lightcone-paper,hqiv-so8-paper};
Tier~1 supplies the variational spine
\cite{hqiv-3d-growth-paper,hqiv-oct-action-paper,hqiv-kirchhoff-paper};
Tier~2 begins with the electroweak and mass ladder
\cite{hqiv-mass-spectrum-paper}. The present paper closes the
loop: every sector of the textbook Standard-Model Lagrangian
$\mathcal{L}_{\rm SM}$ is identified, sector by sector, with
projections of the HQIV \emph{discrete} action already proved
upstream. Tier-0 records carry minted Zenodo DOIs; Tier-1 and
Tier-2 companions cited above are pre-DOI.

\noindent\textbf{Result.}
The page-long $\mathcal{L}_{\rm SM}$ is not assumed. In
machine-checked Lean (zero \texttt{sorry}), we prove a
constructive sector-by-sector bridge: gauge kinetics from the
abelian octonion kinetic density split by the force-sector map;
fermion minimal coupling from the triality-summed source slot;
Higgs kinetic and potential from the octonion-scalar norm and
the lock-in vev; Yukawa weights from the two-step charged-lepton
resonance ladder renormalised against a single $\tau$-Planck
scale on the discrete spine. The headline theorem packages all
five sector equalities simultaneously; a parameter witness
certifies that only the three lattice rationals
$\alpha=3/5$, $\gamma=2/5$, $\alpha_{\rm GUT}=1/42$, the
$\tau$-Planck anchor, and the two resonance steps enter---no
PDG / $\overline{\rm MS}$ mass table, no fitted Yukawa matrix,
no external lattice input. Reader-friendly anchors with
clickable Lean links are in Appendix~\ref{app:lean-index}.

\noindent\textbf{Single active scale witness.}
Mass-sector numerals in the Yukawa block are \emph{outputs} of
the geometric mass functional and resonance ladder under the
same discipline as \cite{hqiv-mass-spectrum-paper,hqiv-kirchhoff-paper}:
exactly one dimensionful witness pins the chart per solve (default:
proton lock-in at $M=\mathrm{referenceM}=4$). PDG centrals and
CODATA couplings are comparison layers, never simultaneous solve
anchors.

\noindent\textbf{Patch-ontology reader contract.}
Per the patch-theory reader contract below, the discrete patch
layer is the ontology. Smooth-manifold textbook notation for
$\mathcal{L}_{\rm SM}$ is a calculation-approximation rewrite
for cross-citation only---not a foundational target, not a
mesh-refinement programme, and not required for validity of the
discrete bridge. At the load-bearing steps (shell index $m$,
resonance ladder, lock-in vev, and the rationals above) a literal
continuum rewrite would erase the discrete index doing the
predictive work.

\noindent\textbf{Status.}
What is proved: the five sector equalities, the total-action
split, the parameter witness, and Yukawa linearity over flavours.
What is named explicitly and \emph{not} claimed: anomaly
cancellation, measure choice, and the full non-abelian holonomy
upgrade (scaffolded upstream but not closed here).
\end{abstract}
\maketitle

% Shared reader contract: patch QFT + discrete numerics.
% From papers/*.tex:     % Shared reader contract: patch QFT + discrete numerics.
% From papers/*.tex:     % Shared reader contract: patch QFT + discrete numerics.
% From papers/*.tex:     \input{include/patch_theory_messaging.tex}
% From papers/paper/*.tex: \input{../include/patch_theory_messaging.tex}

\paragraph{Patch QFT and discrete numerics (reader contract).}
HQIV (Horizon-Quantized Informational Vacuum) is formulated as a \emph{patch theory}: quantum fields, actions, and physical readouts are attached to discrete patches---shell indices $m\in\mathbb{N}$, finite spacetime index types (e.g.\ $\mathrm{Fin}\,4$), and horizon-limited charts---rather than to a hypothetical fundamental smooth spacetime obtained as a mesh-refinement or ultraviolet limit.
Accordingly, when this text refers to ``QFT,'' it means \emph{patch QFT}: patching rules, finite measurement ledgers, and variational identities on the discrete spine; it is not the claim that the theory is \emph{defined} by recovering ordinary continuum quantum field theory through such a limit.
The numbers that admit the sharpest statements---mass and coupling ladders, curvature ratios, shell thermodynamics, and rational witnesses in \texttt{HQIV\_LEAN}---are objects of \emph{discrete mathematics} on $\mathbb{N}$ (combinatorial growth, cumulative simplex ledgers) together with the ordered field $\mathbb{R}$ as used in the proof assistant for analysis, bounds, and phenomenological display.
Smooth-manifold or continuum field notation, when it appears, is a \emph{translation layer} for comparison with standard physics literature; it does not supersede the discrete definitions.

\paragraph{A continuum is not required, and in places would destroy these results.}
The discrete patch layer is \emph{not} a numerical approximation to a more fundamental continuum theory awaiting future construction; it is the ontology.  Where the literature treats ``the continuum limit'' as the goal, HQIV treats it as a \emph{calculation approximation} for comparison with standard formulas.  The continuum form is an optional convenience for comparison; the proved discrete statement is what carries the physics.  In several load-bearing places---the curvature imprint $\delta_E(m)$ at shell $m$, the harmonic ladder masses $m_\tau\to m_\mu\to m_e$, the lock-in vev $v=\sqrt{\eta_{\rm paper}\,\Omega_k(m_{\rm lockin};m_{\rm lockin})}$, the proton anchor at $m=\mathrm{referenceM}=4$, the baryon-asymmetry $\eta\propto 6^7\sqrt{3}/(m+1)$, and the rational coefficients $\alpha=3/5$, $\gamma=2/5$, $\alpha_{\rm GUT}=1/42$---taking a literal $m\to\infty$ or sub-Planck continuum limit would \emph{erase} the discrete index that is doing the predictive work.  Continuum forms of these objects are therefore not preferred refinements; they are degenerate, information-losing readouts of the same patch data.

\paragraph{An exception: $\mathfrak{so}(8)$ as a de-facto continuum algebra from a discrete construction.}
Principled exception to the discrete-only stance: the closed Lie algebra $\mathfrak{so}(8) \supseteq G_2\cup\{\Delta\}$ generated from the Fano octonion left-multiplication table and the phase-lift generator $\Delta\in(\mathrm{e}_1,\mathrm{e}_7)$.  Although $\mathfrak{so}(8)$ is a continuous (28-dimensional real) Lie algebra, it is built here entirely from \emph{discrete} ingredients (the seven imaginary octonion units, a finite Fano table, and one $\pi/2$ rotation), and its closure to 28 dimensions is a finite linear-algebra certificate (see the \texttt{Hqiv.SO8Closure} / \texttt{Hqiv.SO8ClosureInterface} stack).  The Lie-group structure is therefore a \emph{de-facto continuum algebra obtained from a discrete construction}: continuous in parameter space, but with no continuum spacetime, no continuum measure, and no UV limit attached.  When this manuscript or its companions write $\exp(t\,\Delta)\in\mathrm{SO}(8)$ or treat $\mathfrak{so}(8)$ rotations as smooth, those are statements internal to a finite-dimensional Lie group whose generators are certified by discrete data; they are \emph{not} a continuum-spacetime commitment and do not require a continuum-QFT scaffolding.

% From papers/paper/*.tex: % Shared reader contract: patch QFT + discrete numerics.
% From papers/*.tex:     \input{include/patch_theory_messaging.tex}
% From papers/paper/*.tex: \input{../include/patch_theory_messaging.tex}

\paragraph{Patch QFT and discrete numerics (reader contract).}
HQIV (Horizon-Quantized Informational Vacuum) is formulated as a \emph{patch theory}: quantum fields, actions, and physical readouts are attached to discrete patches---shell indices $m\in\mathbb{N}$, finite spacetime index types (e.g.\ $\mathrm{Fin}\,4$), and horizon-limited charts---rather than to a hypothetical fundamental smooth spacetime obtained as a mesh-refinement or ultraviolet limit.
Accordingly, when this text refers to ``QFT,'' it means \emph{patch QFT}: patching rules, finite measurement ledgers, and variational identities on the discrete spine; it is not the claim that the theory is \emph{defined} by recovering ordinary continuum quantum field theory through such a limit.
The numbers that admit the sharpest statements---mass and coupling ladders, curvature ratios, shell thermodynamics, and rational witnesses in \texttt{HQIV\_LEAN}---are objects of \emph{discrete mathematics} on $\mathbb{N}$ (combinatorial growth, cumulative simplex ledgers) together with the ordered field $\mathbb{R}$ as used in the proof assistant for analysis, bounds, and phenomenological display.
Smooth-manifold or continuum field notation, when it appears, is a \emph{translation layer} for comparison with standard physics literature; it does not supersede the discrete definitions.

\paragraph{A continuum is not required, and in places would destroy these results.}
The discrete patch layer is \emph{not} a numerical approximation to a more fundamental continuum theory awaiting future construction; it is the ontology.  Where the literature treats ``the continuum limit'' as the goal, HQIV treats it as a \emph{calculation approximation} for comparison with standard formulas.  The continuum form is an optional convenience for comparison; the proved discrete statement is what carries the physics.  In several load-bearing places---the curvature imprint $\delta_E(m)$ at shell $m$, the harmonic ladder masses $m_\tau\to m_\mu\to m_e$, the lock-in vev $v=\sqrt{\eta_{\rm paper}\,\Omega_k(m_{\rm lockin};m_{\rm lockin})}$, the proton anchor at $m=\mathrm{referenceM}=4$, the baryon-asymmetry $\eta\propto 6^7\sqrt{3}/(m+1)$, and the rational coefficients $\alpha=3/5$, $\gamma=2/5$, $\alpha_{\rm GUT}=1/42$---taking a literal $m\to\infty$ or sub-Planck continuum limit would \emph{erase} the discrete index that is doing the predictive work.  Continuum forms of these objects are therefore not preferred refinements; they are degenerate, information-losing readouts of the same patch data.

\paragraph{An exception: $\mathfrak{so}(8)$ as a de-facto continuum algebra from a discrete construction.}
Principled exception to the discrete-only stance: the closed Lie algebra $\mathfrak{so}(8) \supseteq G_2\cup\{\Delta\}$ generated from the Fano octonion left-multiplication table and the phase-lift generator $\Delta\in(\mathrm{e}_1,\mathrm{e}_7)$.  Although $\mathfrak{so}(8)$ is a continuous (28-dimensional real) Lie algebra, it is built here entirely from \emph{discrete} ingredients (the seven imaginary octonion units, a finite Fano table, and one $\pi/2$ rotation), and its closure to 28 dimensions is a finite linear-algebra certificate (see the \texttt{Hqiv.SO8Closure} / \texttt{Hqiv.SO8ClosureInterface} stack).  The Lie-group structure is therefore a \emph{de-facto continuum algebra obtained from a discrete construction}: continuous in parameter space, but with no continuum spacetime, no continuum measure, and no UV limit attached.  When this manuscript or its companions write $\exp(t\,\Delta)\in\mathrm{SO}(8)$ or treat $\mathfrak{so}(8)$ rotations as smooth, those are statements internal to a finite-dimensional Lie group whose generators are certified by discrete data; they are \emph{not} a continuum-spacetime commitment and do not require a continuum-QFT scaffolding.


\paragraph{Patch QFT and discrete numerics (reader contract).}
HQIV (Horizon-Quantized Informational Vacuum) is formulated as a \emph{patch theory}: quantum fields, actions, and physical readouts are attached to discrete patches---shell indices $m\in\mathbb{N}$, finite spacetime index types (e.g.\ $\mathrm{Fin}\,4$), and horizon-limited charts---rather than to a hypothetical fundamental smooth spacetime obtained as a mesh-refinement or ultraviolet limit.
Accordingly, when this text refers to ``QFT,'' it means \emph{patch QFT}: patching rules, finite measurement ledgers, and variational identities on the discrete spine; it is not the claim that the theory is \emph{defined} by recovering ordinary continuum quantum field theory through such a limit.
The numbers that admit the sharpest statements---mass and coupling ladders, curvature ratios, shell thermodynamics, and rational witnesses in \texttt{HQIV\_LEAN}---are objects of \emph{discrete mathematics} on $\mathbb{N}$ (combinatorial growth, cumulative simplex ledgers) together with the ordered field $\mathbb{R}$ as used in the proof assistant for analysis, bounds, and phenomenological display.
Smooth-manifold or continuum field notation, when it appears, is a \emph{translation layer} for comparison with standard physics literature; it does not supersede the discrete definitions.

\paragraph{A continuum is not required, and in places would destroy these results.}
The discrete patch layer is \emph{not} a numerical approximation to a more fundamental continuum theory awaiting future construction; it is the ontology.  Where the literature treats ``the continuum limit'' as the goal, HQIV treats it as a \emph{calculation approximation} for comparison with standard formulas.  The continuum form is an optional convenience for comparison; the proved discrete statement is what carries the physics.  In several load-bearing places---the curvature imprint $\delta_E(m)$ at shell $m$, the harmonic ladder masses $m_\tau\to m_\mu\to m_e$, the lock-in vev $v=\sqrt{\eta_{\rm paper}\,\Omega_k(m_{\rm lockin};m_{\rm lockin})}$, the proton anchor at $m=\mathrm{referenceM}=4$, the baryon-asymmetry $\eta\propto 6^7\sqrt{3}/(m+1)$, and the rational coefficients $\alpha=3/5$, $\gamma=2/5$, $\alpha_{\rm GUT}=1/42$---taking a literal $m\to\infty$ or sub-Planck continuum limit would \emph{erase} the discrete index that is doing the predictive work.  Continuum forms of these objects are therefore not preferred refinements; they are degenerate, information-losing readouts of the same patch data.

\paragraph{An exception: $\mathfrak{so}(8)$ as a de-facto continuum algebra from a discrete construction.}
Principled exception to the discrete-only stance: the closed Lie algebra $\mathfrak{so}(8) \supseteq G_2\cup\{\Delta\}$ generated from the Fano octonion left-multiplication table and the phase-lift generator $\Delta\in(\mathrm{e}_1,\mathrm{e}_7)$.  Although $\mathfrak{so}(8)$ is a continuous (28-dimensional real) Lie algebra, it is built here entirely from \emph{discrete} ingredients (the seven imaginary octonion units, a finite Fano table, and one $\pi/2$ rotation), and its closure to 28 dimensions is a finite linear-algebra certificate (see the \texttt{Hqiv.SO8Closure} / \texttt{Hqiv.SO8ClosureInterface} stack).  The Lie-group structure is therefore a \emph{de-facto continuum algebra obtained from a discrete construction}: continuous in parameter space, but with no continuum spacetime, no continuum measure, and no UV limit attached.  When this manuscript or its companions write $\exp(t\,\Delta)\in\mathrm{SO}(8)$ or treat $\mathfrak{so}(8)$ rotations as smooth, those are statements internal to a finite-dimensional Lie group whose generators are certified by discrete data; they are \emph{not} a continuum-spacetime commitment and do not require a continuum-QFT scaffolding.

% From papers/paper/*.tex: % Shared reader contract: patch QFT + discrete numerics.
% From papers/*.tex:     % Shared reader contract: patch QFT + discrete numerics.
% From papers/*.tex:     \input{include/patch_theory_messaging.tex}
% From papers/paper/*.tex: \input{../include/patch_theory_messaging.tex}

\paragraph{Patch QFT and discrete numerics (reader contract).}
HQIV (Horizon-Quantized Informational Vacuum) is formulated as a \emph{patch theory}: quantum fields, actions, and physical readouts are attached to discrete patches---shell indices $m\in\mathbb{N}$, finite spacetime index types (e.g.\ $\mathrm{Fin}\,4$), and horizon-limited charts---rather than to a hypothetical fundamental smooth spacetime obtained as a mesh-refinement or ultraviolet limit.
Accordingly, when this text refers to ``QFT,'' it means \emph{patch QFT}: patching rules, finite measurement ledgers, and variational identities on the discrete spine; it is not the claim that the theory is \emph{defined} by recovering ordinary continuum quantum field theory through such a limit.
The numbers that admit the sharpest statements---mass and coupling ladders, curvature ratios, shell thermodynamics, and rational witnesses in \texttt{HQIV\_LEAN}---are objects of \emph{discrete mathematics} on $\mathbb{N}$ (combinatorial growth, cumulative simplex ledgers) together with the ordered field $\mathbb{R}$ as used in the proof assistant for analysis, bounds, and phenomenological display.
Smooth-manifold or continuum field notation, when it appears, is a \emph{translation layer} for comparison with standard physics literature; it does not supersede the discrete definitions.

\paragraph{A continuum is not required, and in places would destroy these results.}
The discrete patch layer is \emph{not} a numerical approximation to a more fundamental continuum theory awaiting future construction; it is the ontology.  Where the literature treats ``the continuum limit'' as the goal, HQIV treats it as a \emph{calculation approximation} for comparison with standard formulas.  The continuum form is an optional convenience for comparison; the proved discrete statement is what carries the physics.  In several load-bearing places---the curvature imprint $\delta_E(m)$ at shell $m$, the harmonic ladder masses $m_\tau\to m_\mu\to m_e$, the lock-in vev $v=\sqrt{\eta_{\rm paper}\,\Omega_k(m_{\rm lockin};m_{\rm lockin})}$, the proton anchor at $m=\mathrm{referenceM}=4$, the baryon-asymmetry $\eta\propto 6^7\sqrt{3}/(m+1)$, and the rational coefficients $\alpha=3/5$, $\gamma=2/5$, $\alpha_{\rm GUT}=1/42$---taking a literal $m\to\infty$ or sub-Planck continuum limit would \emph{erase} the discrete index that is doing the predictive work.  Continuum forms of these objects are therefore not preferred refinements; they are degenerate, information-losing readouts of the same patch data.

\paragraph{An exception: $\mathfrak{so}(8)$ as a de-facto continuum algebra from a discrete construction.}
Principled exception to the discrete-only stance: the closed Lie algebra $\mathfrak{so}(8) \supseteq G_2\cup\{\Delta\}$ generated from the Fano octonion left-multiplication table and the phase-lift generator $\Delta\in(\mathrm{e}_1,\mathrm{e}_7)$.  Although $\mathfrak{so}(8)$ is a continuous (28-dimensional real) Lie algebra, it is built here entirely from \emph{discrete} ingredients (the seven imaginary octonion units, a finite Fano table, and one $\pi/2$ rotation), and its closure to 28 dimensions is a finite linear-algebra certificate (see the \texttt{Hqiv.SO8Closure} / \texttt{Hqiv.SO8ClosureInterface} stack).  The Lie-group structure is therefore a \emph{de-facto continuum algebra obtained from a discrete construction}: continuous in parameter space, but with no continuum spacetime, no continuum measure, and no UV limit attached.  When this manuscript or its companions write $\exp(t\,\Delta)\in\mathrm{SO}(8)$ or treat $\mathfrak{so}(8)$ rotations as smooth, those are statements internal to a finite-dimensional Lie group whose generators are certified by discrete data; they are \emph{not} a continuum-spacetime commitment and do not require a continuum-QFT scaffolding.

% From papers/paper/*.tex: % Shared reader contract: patch QFT + discrete numerics.
% From papers/*.tex:     \input{include/patch_theory_messaging.tex}
% From papers/paper/*.tex: \input{../include/patch_theory_messaging.tex}

\paragraph{Patch QFT and discrete numerics (reader contract).}
HQIV (Horizon-Quantized Informational Vacuum) is formulated as a \emph{patch theory}: quantum fields, actions, and physical readouts are attached to discrete patches---shell indices $m\in\mathbb{N}$, finite spacetime index types (e.g.\ $\mathrm{Fin}\,4$), and horizon-limited charts---rather than to a hypothetical fundamental smooth spacetime obtained as a mesh-refinement or ultraviolet limit.
Accordingly, when this text refers to ``QFT,'' it means \emph{patch QFT}: patching rules, finite measurement ledgers, and variational identities on the discrete spine; it is not the claim that the theory is \emph{defined} by recovering ordinary continuum quantum field theory through such a limit.
The numbers that admit the sharpest statements---mass and coupling ladders, curvature ratios, shell thermodynamics, and rational witnesses in \texttt{HQIV\_LEAN}---are objects of \emph{discrete mathematics} on $\mathbb{N}$ (combinatorial growth, cumulative simplex ledgers) together with the ordered field $\mathbb{R}$ as used in the proof assistant for analysis, bounds, and phenomenological display.
Smooth-manifold or continuum field notation, when it appears, is a \emph{translation layer} for comparison with standard physics literature; it does not supersede the discrete definitions.

\paragraph{A continuum is not required, and in places would destroy these results.}
The discrete patch layer is \emph{not} a numerical approximation to a more fundamental continuum theory awaiting future construction; it is the ontology.  Where the literature treats ``the continuum limit'' as the goal, HQIV treats it as a \emph{calculation approximation} for comparison with standard formulas.  The continuum form is an optional convenience for comparison; the proved discrete statement is what carries the physics.  In several load-bearing places---the curvature imprint $\delta_E(m)$ at shell $m$, the harmonic ladder masses $m_\tau\to m_\mu\to m_e$, the lock-in vev $v=\sqrt{\eta_{\rm paper}\,\Omega_k(m_{\rm lockin};m_{\rm lockin})}$, the proton anchor at $m=\mathrm{referenceM}=4$, the baryon-asymmetry $\eta\propto 6^7\sqrt{3}/(m+1)$, and the rational coefficients $\alpha=3/5$, $\gamma=2/5$, $\alpha_{\rm GUT}=1/42$---taking a literal $m\to\infty$ or sub-Planck continuum limit would \emph{erase} the discrete index that is doing the predictive work.  Continuum forms of these objects are therefore not preferred refinements; they are degenerate, information-losing readouts of the same patch data.

\paragraph{An exception: $\mathfrak{so}(8)$ as a de-facto continuum algebra from a discrete construction.}
Principled exception to the discrete-only stance: the closed Lie algebra $\mathfrak{so}(8) \supseteq G_2\cup\{\Delta\}$ generated from the Fano octonion left-multiplication table and the phase-lift generator $\Delta\in(\mathrm{e}_1,\mathrm{e}_7)$.  Although $\mathfrak{so}(8)$ is a continuous (28-dimensional real) Lie algebra, it is built here entirely from \emph{discrete} ingredients (the seven imaginary octonion units, a finite Fano table, and one $\pi/2$ rotation), and its closure to 28 dimensions is a finite linear-algebra certificate (see the \texttt{Hqiv.SO8Closure} / \texttt{Hqiv.SO8ClosureInterface} stack).  The Lie-group structure is therefore a \emph{de-facto continuum algebra obtained from a discrete construction}: continuous in parameter space, but with no continuum spacetime, no continuum measure, and no UV limit attached.  When this manuscript or its companions write $\exp(t\,\Delta)\in\mathrm{SO}(8)$ or treat $\mathfrak{so}(8)$ rotations as smooth, those are statements internal to a finite-dimensional Lie group whose generators are certified by discrete data; they are \emph{not} a continuum-spacetime commitment and do not require a continuum-QFT scaffolding.


\paragraph{Patch QFT and discrete numerics (reader contract).}
HQIV (Horizon-Quantized Informational Vacuum) is formulated as a \emph{patch theory}: quantum fields, actions, and physical readouts are attached to discrete patches---shell indices $m\in\mathbb{N}$, finite spacetime index types (e.g.\ $\mathrm{Fin}\,4$), and horizon-limited charts---rather than to a hypothetical fundamental smooth spacetime obtained as a mesh-refinement or ultraviolet limit.
Accordingly, when this text refers to ``QFT,'' it means \emph{patch QFT}: patching rules, finite measurement ledgers, and variational identities on the discrete spine; it is not the claim that the theory is \emph{defined} by recovering ordinary continuum quantum field theory through such a limit.
The numbers that admit the sharpest statements---mass and coupling ladders, curvature ratios, shell thermodynamics, and rational witnesses in \texttt{HQIV\_LEAN}---are objects of \emph{discrete mathematics} on $\mathbb{N}$ (combinatorial growth, cumulative simplex ledgers) together with the ordered field $\mathbb{R}$ as used in the proof assistant for analysis, bounds, and phenomenological display.
Smooth-manifold or continuum field notation, when it appears, is a \emph{translation layer} for comparison with standard physics literature; it does not supersede the discrete definitions.

\paragraph{A continuum is not required, and in places would destroy these results.}
The discrete patch layer is \emph{not} a numerical approximation to a more fundamental continuum theory awaiting future construction; it is the ontology.  Where the literature treats ``the continuum limit'' as the goal, HQIV treats it as a \emph{calculation approximation} for comparison with standard formulas.  The continuum form is an optional convenience for comparison; the proved discrete statement is what carries the physics.  In several load-bearing places---the curvature imprint $\delta_E(m)$ at shell $m$, the harmonic ladder masses $m_\tau\to m_\mu\to m_e$, the lock-in vev $v=\sqrt{\eta_{\rm paper}\,\Omega_k(m_{\rm lockin};m_{\rm lockin})}$, the proton anchor at $m=\mathrm{referenceM}=4$, the baryon-asymmetry $\eta\propto 6^7\sqrt{3}/(m+1)$, and the rational coefficients $\alpha=3/5$, $\gamma=2/5$, $\alpha_{\rm GUT}=1/42$---taking a literal $m\to\infty$ or sub-Planck continuum limit would \emph{erase} the discrete index that is doing the predictive work.  Continuum forms of these objects are therefore not preferred refinements; they are degenerate, information-losing readouts of the same patch data.

\paragraph{An exception: $\mathfrak{so}(8)$ as a de-facto continuum algebra from a discrete construction.}
Principled exception to the discrete-only stance: the closed Lie algebra $\mathfrak{so}(8) \supseteq G_2\cup\{\Delta\}$ generated from the Fano octonion left-multiplication table and the phase-lift generator $\Delta\in(\mathrm{e}_1,\mathrm{e}_7)$.  Although $\mathfrak{so}(8)$ is a continuous (28-dimensional real) Lie algebra, it is built here entirely from \emph{discrete} ingredients (the seven imaginary octonion units, a finite Fano table, and one $\pi/2$ rotation), and its closure to 28 dimensions is a finite linear-algebra certificate (see the \texttt{Hqiv.SO8Closure} / \texttt{Hqiv.SO8ClosureInterface} stack).  The Lie-group structure is therefore a \emph{de-facto continuum algebra obtained from a discrete construction}: continuous in parameter space, but with no continuum spacetime, no continuum measure, and no UV limit attached.  When this manuscript or its companions write $\exp(t\,\Delta)\in\mathrm{SO}(8)$ or treat $\mathfrak{so}(8)$ rotations as smooth, those are statements internal to a finite-dimensional Lie group whose generators are certified by discrete data; they are \emph{not} a continuum-spacetime commitment and do not require a continuum-QFT scaffolding.


\paragraph{Patch QFT and discrete numerics (reader contract).}
HQIV (Horizon-Quantized Informational Vacuum) is formulated as a \emph{patch theory}: quantum fields, actions, and physical readouts are attached to discrete patches---shell indices $m\in\mathbb{N}$, finite spacetime index types (e.g.\ $\mathrm{Fin}\,4$), and horizon-limited charts---rather than to a hypothetical fundamental smooth spacetime obtained as a mesh-refinement or ultraviolet limit.
Accordingly, when this text refers to ``QFT,'' it means \emph{patch QFT}: patching rules, finite measurement ledgers, and variational identities on the discrete spine; it is not the claim that the theory is \emph{defined} by recovering ordinary continuum quantum field theory through such a limit.
The numbers that admit the sharpest statements---mass and coupling ladders, curvature ratios, shell thermodynamics, and rational witnesses in \texttt{HQIV\_LEAN}---are objects of \emph{discrete mathematics} on $\mathbb{N}$ (combinatorial growth, cumulative simplex ledgers) together with the ordered field $\mathbb{R}$ as used in the proof assistant for analysis, bounds, and phenomenological display.
Smooth-manifold or continuum field notation, when it appears, is a \emph{translation layer} for comparison with standard physics literature; it does not supersede the discrete definitions.

\paragraph{A continuum is not required, and in places would destroy these results.}
The discrete patch layer is \emph{not} a numerical approximation to a more fundamental continuum theory awaiting future construction; it is the ontology.  Where the literature treats ``the continuum limit'' as the goal, HQIV treats it as a \emph{calculation approximation} for comparison with standard formulas.  The continuum form is an optional convenience for comparison; the proved discrete statement is what carries the physics.  In several load-bearing places---the curvature imprint $\delta_E(m)$ at shell $m$, the harmonic ladder masses $m_\tau\to m_\mu\to m_e$, the lock-in vev $v=\sqrt{\eta_{\rm paper}\,\Omega_k(m_{\rm lockin};m_{\rm lockin})}$, the proton anchor at $m=\mathrm{referenceM}=4$, the baryon-asymmetry $\eta\propto 6^7\sqrt{3}/(m+1)$, and the rational coefficients $\alpha=3/5$, $\gamma=2/5$, $\alpha_{\rm GUT}=1/42$---taking a literal $m\to\infty$ or sub-Planck continuum limit would \emph{erase} the discrete index that is doing the predictive work.  Continuum forms of these objects are therefore not preferred refinements; they are degenerate, information-losing readouts of the same patch data.

\paragraph{An exception: $\mathfrak{so}(8)$ as a de-facto continuum algebra from a discrete construction.}
Principled exception to the discrete-only stance: the closed Lie algebra $\mathfrak{so}(8) \supseteq G_2\cup\{\Delta\}$ generated from the Fano octonion left-multiplication table and the phase-lift generator $\Delta\in(\mathrm{e}_1,\mathrm{e}_7)$.  Although $\mathfrak{so}(8)$ is a continuous (28-dimensional real) Lie algebra, it is built here entirely from \emph{discrete} ingredients (the seven imaginary octonion units, a finite Fano table, and one $\pi/2$ rotation), and its closure to 28 dimensions is a finite linear-algebra certificate (see the \texttt{Hqiv.SO8Closure} / \texttt{Hqiv.SO8ClosureInterface} stack).  The Lie-group structure is therefore a \emph{de-facto continuum algebra obtained from a discrete construction}: continuous in parameter space, but with no continuum spacetime, no continuum measure, and no UV limit attached.  When this manuscript or its companions write $\exp(t\,\Delta)\in\mathrm{SO}(8)$ or treat $\mathfrak{so}(8)$ rotations as smooth, those are statements internal to a finite-dimensional Lie group whose generators are certified by discrete data; they are \emph{not} a continuum-spacetime commitment and do not require a continuum-QFT scaffolding.


\section{Scope and claim discipline}
\label{sec:scope}

\subsection{Predecessor: discrete O--Maxwell action}

\cite{hqiv-oct-action-paper} packages the discrete octonion-indexed gauge
potential $A^a{}_\mu$ (channels $a\in\mathrm{Fin}\,8$, spacetime indices
$\mu,\nu\in\mathrm{Fin}\,4$), field strength
$F^a{}_{\mu\nu}(A)=A^a{}_\nu - A^a{}_\mu$ (the
\hyperref[lean:F-from-A]{field-strength constructor}), the kinetic density
$L_O^{\rm kin}(A)$ (the
\hyperref[lean:LO-kin]{octonion kinetic density}), the source interaction
$L_O^{\rm src}(J_{\rm src},A)$ (the
\hyperref[lean:LO-src]{octonion source coupling}), and the $\phi$-coupling
slot $L_O^\phi(A,\phi_{\rm val})$ (the
\hyperref[lean:LO-phi]{scalar $\phi$-coupling slot}) into the single
Lagrangian density (the
\hyperref[lean:LO-maxwell]{general O-Maxwell Lagrangian density})
\[
\mathcal{L}_O[J_{\rm src}](A,\phi_{\rm val})
= L_O^{\rm kin}(A) + 4\pi\,L_O^{\rm src}(J_{\rm src},A)
+ L_O^\phi(A,\phi_{\rm val}).
\]
The gravitational constraint scalar $S_{\rm HQVM\,grav}$ (the
\hyperref[lean:SHQVM-grav]{HQVM gravitational constraint}) vanishes iff the
HQVM Friedmann equation holds (the
\hyperref[lean:friedmann-iff]{Friedmann equivalence theorem}). The total cell
(the \hyperref[lean:action-total]{HQIV total action}) is
$\mathcal{L}_O+S_{\rm HQVM\,grav}$. \cite{hqiv-mass-spectrum-paper} adds the
charged-lepton resonance ladder (the
\hyperref[lean:resonance-steps]{two resonance steps}), the $\tau$-Planck
anchor (the \hyperref[lean:m-tau-pl]{tau Planck mass}), and the
generation-level mass functional (the
\hyperref[lean:sm-mass-label]{geometric SM mass functional}) that ties an SM
flavour label to a Planck-unit rest mass with no PDG input.
\cite{hqiv-so8-paper} and the
\hyperref[lean:sm-embedding]{SM embedding module} fix the
$\mathrm{SU}(2)_L$ and hypercharge generators inside the certified
$G_2\cup\{\Delta\}\Rightarrow\mathfrak{so}(8)$ closure, exhaust the
standard branching rules on $8_s$, and certify the three triality irreps
that we read as generation labels.

\subsection{What this paper adds}
\label{subsec:adds}

The new Lean module is the
\hyperref[lean:sm-lagrangian-module]{Standard-Model Lagrangian bridge}
(zero \texttt{sorry}, builds clean under \texttt{lake build}); its principal
contributions are:

\begin{itemize}[leftmargin=2em]
\item \textbf{Sector splitting of the abelian octonion kinetic}
  (Sect.~\ref{sec:gauge}): the
  \hyperref[lean:LO-kinetic-split]{three-sector kinetic split} writes
  $L_O^{\rm kin}=L_{O,{\rm EM}}^{\rm kin}+L_{O,{\rm Weak}}^{\rm kin}+L_{O,{\rm Strong}}^{\rm kin}$
  under the explicit force-sector assignment (the
  \hyperref[lean:O-to-sector]{force-sector map}) from the
  \hyperref[lean:forces]{Forces module}; sector aggregates project to
  the textbook $B$, $W$, $G$ kinetic blocks; the SM gauge total equals
  $L_O^{\rm kin}$ exactly (the
  \hyperref[lean:gauge-eq-kin]{gauge--kinetic equality}).
\item \textbf{Triality-summed fermion coupling} (Sect.~\ref{sec:fermion}):
  per-generation fermion currents on the $8_s$ carrier are summed by two
  applications of the proved $J$-linearity (the
  \hyperref[lean:J-add]{source-coupling additivity lemma}); the
  \hyperref[lean:three-gen]{three-generation equality} packages the result.
\item \textbf{Higgs as the octonion scalar} (Sect.~\ref{sec:higgs}):
  the Higgs kinetic block is the per-direction sum of the
  \hyperref[lean:scalar-norm]{octonion scalar norm}; the
  symmetry-breaking potential is the existing
  \hyperref[lean:higgs-pot]{Higgs potential} with
  $\lambda\mapsto\lambda_{\rm eff}$ (the
  \hyperref[lean:lambda-eff]{effective quartic}) and
  $v\mapsto\texttt{lockinVev}$ (the
  \hyperref[lean:lockin-vev]{lock-in vev}), expanded to the textbook
  quartic form (the
  \hyperref[lean:higgs-expanded]{expanded Higgs potential identity}).
\item \textbf{Yukawa from the resonance ladder} (Sect.~\ref{sec:yukawa}):
  $y_f := \sqrt{2}\,m_f/v$ with $m_f$ from the geometric mass functional;
  the identity $y_f v / \sqrt{2}=m_f$ is proved (the
  \hyperref[lean:yukawa-mass]{Yukawa--mass identity}), and the explicit
  witness exhibits the entire Yukawa block as the two-step ladder (the
  \hyperref[lean:yukawa-resonance]{Yukawa resonance witness}).
\item \textbf{Assembled SM Lagrangian} (Sect.~\ref{sec:assembly}):
  the \hyperref[lean:SM-record]{SM Lagrangian record} and its
  \hyperref[lean:SM-fromHQIV]{from-HQIV constructor} package the five
  sector pieces; the headline theorem (the
  \hyperref[lean:headline]{SM-from-HQIV discrete-action theorem})
  certifies the five sector equalities simultaneously; the total-action
  equality (the
  \hyperref[lean:total-split]{total-action split theorem}) exhibits the
  SM gauge $+$ fermion $+$ $\phi$-coupling $+$ Friedmann split of the HQIV
  total cell when $J_{\rm src}=\sum_{\rm gen} J(\rm gen)$.
\item \textbf{Parameter witness} (Sect.~\ref{sec:parameters}):
  the \hyperref[lean:param-witness]{parameter witness} certifies the three
  lattice rationals $\alpha=3/5$, $\gamma=2/5$, $\alpha_{\rm GUT}=1/42$
  (the \hyperref[lean:alpha-35]{lattice $\alpha$ identity}, the
  \hyperref[lean:gamma-25]{horizon-split $\gamma$ identity}, and the
  \hyperref[lean:alpha-gut]{GUT coupling identity}); together with
  the $\tau$-Planck anchor and the two resonance steps these are the
  entire discrete parameter input to $\mathcal{L}_{\rm SM}$---not a
  simultaneous PDG fit set.
\end{itemize}

\subsection{What is and isn't proved}
\label{subsec:not-claimed}

This module is a \emph{structural bridge}, not a re-derivation of textbook
field equations:

\begin{itemize}[leftmargin=2em]
\item \textbf{Proved (Lean).} The five sector equalities of the
  \hyperref[lean:headline]{SM-from-HQIV discrete-action theorem}; the
  total-action split of the
  \hyperref[lean:total-split]{total-action split theorem}; the
  \hyperref[lean:param-witness]{parameter witness}; linearity of the
  Yukawa sum (the \hyperref[lean:yukawa-sum]{Yukawa sum identity}).
\item \textbf{Black-box ingredients.} The HQIV kinetic, source,
  $\phi$-coupling, and gravitational pieces, together with the Higgs
  potential, lock-in vev, geometric mass functional, effective quartic,
  and force-sector map, are imported from upstream modules; all
  \texttt{sm\_*\_from\_HQIV} theorems are tautological equalities.
  Any smooth-manifold translation is a comparison layer only. Load-bearing
  predictions of \cite{hqiv-oct-action-paper} live on the discrete spine.
\item \textbf{Symbolic kinetic terms.} Fermion and Higgs kinetic shadows
  remain on $\Phi:\mathrm{Fin}\,8\to\mathbb{R}$; no BRST quantisation here.
\item \textbf{Non-abelian holonomy.} Per-sector kinetic uses abelian
  $F^a{}_{\mu\nu}=A^a{}_\nu-A^a{}_\mu$; the weak commutator scaffold is
  upstream and not closed in this module.
\item \textbf{Single witness, no PDG fit.} Yukawa weights come from the
  resonance ladder and $\tau$-Planck anchor on the discrete spine; the Higgs
  vev is the lock-in $v$ above. Under the default proton lock-in witness of
  \cite{hqiv-mass-spectrum-paper,hqiv-kirchhoff-paper}, PDG masses are
  comparison targets, not simultaneous solve anchors.
\end{itemize}

\section{Sector projection of the gauge kinetic block}
\label{sec:gauge}

The HQIV abelian octonion kinetic density (the
\hyperref[lean:LO-kin]{octonion kinetic density} of
\cite{hqiv-oct-action-paper}),
\[
L_O^{\rm kin}(A)
= -\tfrac14\,\sum_{a\in\mathrm{Fin}\,8}\,\sum_{\mu,\nu\in\mathrm{Fin}\,4}
\tfrac{(F^a{}_{\mu\nu}(A))^2}{2},
\]
splits into three sector pieces. The force-sector assignment is the
$O$-component map $\mathrm{forceSector}(a)$ with
$\mathrm{forceSector}(0)=\mathrm{EM}$,
$\mathrm{forceSector}(a)=\mathrm{Weak}$ for $a\in\{1,2,3\}$, and
$\mathrm{forceSector}(a)=\mathrm{Strong}$ for $a\in\{4,5,6,7\}$ (the
\hyperref[lean:O-to-sector]{force-sector map}), fixed in the
\hyperref[lean:forces]{Forces module}; component~$0$ in the quaternionic
subalgebra reduces to classic Maxwell (the
\hyperref[lean:O-zero-EM]{zeroth-component is EM lemma}), the three
remaining quaternionic components carry the weak-like sector, and the four
octonion-extension components carry the colour-like sector.

\begin{definition}[Sector projections of $L_O^{\rm kin}$]
\label{def:LSM-kin}
For $s\in\{\text{EM},\text{Weak},\text{Strong}\}$ and a potential $A$,
\[
L^{\rm kin}_{O,s}(A)
:= -\tfrac14\,\sum_{a\in\mathrm{Fin}\,8} \chi_s(a)\,
\sum_{\mu,\nu\in\mathrm{Fin}\,4}\tfrac{(F^a{}_{\mu\nu}(A))^2}{2},
\]
where $\chi_s(a)=1$ if $\mathrm{forceSector}(a)=s$ and zero otherwise (the
\hyperref[lean:LO-kin-sector]{sector-projected kinetic density}). Aggregating:
\[
\mathcal{L}_B(A) := L^{\rm kin}_{O,{\rm EM}}(A),\qquad
\mathcal{L}_W(A) := L^{\rm kin}_{O,{\rm Weak}}(A),\qquad
\mathcal{L}_G(A) := L^{\rm kin}_{O,{\rm Strong}}(A)
\]
(the \hyperref[lean:LSM-B-kin]{hypercharge kinetic},
\hyperref[lean:LSM-W-kin]{weak kinetic}, and
\hyperref[lean:LSM-G-kin]{strong kinetic} aggregates).
\end{definition}

\begin{theorem}[Three-sector splitting, Lean]
\label{thm:sector-split}
For every $A$,
\[
L_O^{\rm kin}(A)
= L^{\rm kin}_{O,{\rm EM}}(A) + L^{\rm kin}_{O,{\rm Weak}}(A) + L^{\rm kin}_{O,{\rm Strong}}(A) .
\]
This is the \hyperref[lean:LO-kinetic-split]{three-sector kinetic split}.
\end{theorem}

\begin{proof}[Sketch (Lean tactic translation)]
For each octonion channel $a\in\mathrm{Fin}\,8$, exactly one of the three
indicator functions $\chi_{\rm EM}(a), \chi_{\rm Weak}(a), \chi_{\rm Strong}(a)$
is one and the other two vanish (case split on $\mathrm{forceSector}(a)$).
Summing the channel-wise identity over $a$ and applying $-\tfrac14\cdot$
yields the claim. Lean discharges the finite case split with
\texttt{cases}~$+$~\texttt{simp} and the global sum with
\texttt{Finset.sum\_add\_distrib}.
\end{proof}

\begin{corollary}[SM gauge kinetic equals HQIV octonion kinetic]
\label{cor:gauge-aggregate}
$\mathcal{L}_B(A)+\mathcal{L}_W(A)+\mathcal{L}_G(A)=L_O^{\rm kin}(A)$
(the \hyperref[lean:gauge-eq-kin]{gauge--kinetic equality}).
\end{corollary}

\paragraph{Reading.}
The textbook SM gauge kinetic block
$\mathcal{L}_{\rm gauge}=-\tfrac14(B_{\mu\nu}B^{\mu\nu}+W^I_{\mu\nu}W^{I\mu\nu}
+G^a_{\mu\nu}G^{a\mu\nu})$ is the sector-projected aggregate of the single
HQIV object $L_O^{\rm kin}$ when the force-sector map is the identification
fixed by the algebra (colour-preferred axis and hypercharge block of the
\hyperref[lean:sm-embedding]{SM embedding module}). No new coupling, no new
field strength: the gauge-kinetic page of $\mathcal{L}_{\rm SM}$ is one
discrete object with a sector partition.

\section{Fermion sector: minimal coupling and three generations}
\label{sec:fermion}

The textbook block $i\bar\psi\gamma^\mu D_\mu\psi$ summed over one chiral
SM generation is supported on the 8-dimensional left-handed $8_s$ carrier of
the \hyperref[lean:sm-embedding]{SM embedding module}, plus 8 right-handed
components in $8_c$ (see the branching rule on $8_s$ documented there).
At the abelian level needed for the present bridge, the Dirac bilinear is
encoded as a current $J:\mathrm{Fin}\,8\to\mathrm{Fin}\,4\to\mathbb{R}$ on
the 8s carrier, and the minimal coupling reads off the proved $J\cdot A$ slot
$L_O^{\rm src}$.

\begin{definition}[Per-generation minimal coupling]
\label{def:fermion-coupling}
For a triality generation label $\mathrm{gen}\in\texttt{So8RepIndex}$ and a
generation current $J_{\rm gen}$,
\[
L_{\rm fermion}^{\rm SM}(\mathrm{gen}, J_{\rm gen}, A) := L_O^{\rm src}(J_{\rm gen},A)
\]
(the \hyperref[lean:fermion-minimal]{per-generation minimal coupling}).
The three SM generations correspond to the three Spin(8) irreps
$8_v, 8_{s^+}, 8_{s^-}$ counted by the
\hyperref[lean:three-gen-reps]{three-generations-from-triality theorem}
($\mathrm{card}\,\texttt{So8RepIndex}=3$).
\end{definition}

\begin{theorem}[Triality-summed fermion coupling, Lean]
\label{thm:three-generations}
For any triality-indexed family of currents $J:\texttt{So8RepIndex}\to(\mathrm{Fin}\,8\to\mathrm{Fin}\,4\to\mathbb{R})$,
\[
\sum_{\mathrm{gen}\in\{8_v,8_{s^+},8_{s^-}\}} L_{\rm fermion}^{\rm SM}(\mathrm{gen},J(\mathrm{gen}),A)
= L_O^{\rm src}\!\left(\sum_{\rm gen} J(\rm gen),\;A\right).
\]
This is the \hyperref[lean:three-gen]{three-generation equality}; its proof
is two applications of the
\hyperref[lean:J-add]{source-coupling additivity lemma}
(Lemma~3.1 of \cite{hqiv-oct-action-paper}).
\end{theorem}

\begin{remark}[$8_v$, $8_{s^+}$, $8_{s^-}$ vs.\ generation labelling]
The \hyperref[lean:sm-gr]{SM--GR unification module} assigns generation
$1\to 8_v$, generation $2\to 8_{s^+}$, generation $3\to 8_{s^-}$ via the
\hyperref[lean:sm-gen-index]{generation index map}; this is a
representation-counting fact, not a mass theorem. The mass hierarchy itself
is read off the \hyperref[lean:lepton-resonance]{charged-lepton resonance
module} via the \hyperref[lean:resonance-product]{resonance-product factor}
(Section~\ref{sec:yukawa} below).
\end{remark}

\paragraph{Reading.}
The textbook fermion sector
$\mathcal{L}_{\rm fermion}=\sum_{\rm gen}i\bar\psi_{\rm gen}\gamma^\mu D_\mu\psi_{\rm gen}$
is, at the bookkeeping level relevant to the discrete action, the abelian
source coupling for the sum of the three triality-indexed currents on
the $8_s$ carrier. The covariant derivative $D_\mu$ encodes the same gauge
connection $A$ that supplies the sector kinetic block of
Section~\ref{sec:gauge}; no additional gauge structure is needed.

\section{Higgs sector: octonion scalar, vev, and quartic}
\label{sec:higgs}

The HQIV scalar lives on the same $\mathrm{Fin}\,8\to\mathbb{R}$ carrier as
the gauge potential (the \hyperref[lean:weak-higgs]{weak/Higgs scaffold});
the Higgs kinetic shadow of the textbook $(D_\mu\Phi)^\dagger(D^\mu\Phi)$ is
read off the \hyperref[lean:scalar-norm]{octonion scalar norm}, summed over
spacetime directions; the symmetry-breaking potential is the
\hyperref[lean:higgs-pot]{Higgs potential} at the calibrated
\hyperref[lean:lockin-vev]{lock-in vev}.

\begin{definition}[Higgs sector projection]
\label{def:higgs}
For a per-direction scalar slot $\Phi_{\rm dir}:\mathrm{Fin}\,4\to(\mathrm{Fin}\,8\to\mathbb{R})$,
\[
\mathcal{L}_{\rm Higgs}^{\rm kin}(\Phi_{\rm dir})
:= \sum_{\mu\in\mathrm{Fin}\,4}\,\mathrm{scalarNormSq}(\Phi_{\rm dir}(\mu))
\]
(the \hyperref[lean:LSM-Higgs-kin]{SM Higgs kinetic shadow}). For a scalar
configuration $\Phi$,
\[
\mathcal{L}_{\rm Higgs}^{\rm pot}(\Phi)
:= \texttt{higgsPotential}(\lambda_{\rm eff},\,\texttt{lockinVev},\,\Phi)
= \lambda_{\rm eff}\bigl(|\Phi|^2 - \texttt{lockinVev}^2\bigr)^2
\]
(the \hyperref[lean:LSM-Higgs-pot]{SM Higgs potential shadow}), with
$\lambda_{\rm eff}$ from the \hyperref[lean:sm-gr]{SM--GR unification module}
and the lock-in vev from the weak/Higgs scaffold,
\[
\texttt{lockinVev}
:= \sqrt{\eta_{\rm paper}\,\Omega_k(m_{\rm lockin};m_{\rm lockin})}.
\]
\end{definition}

\begin{theorem}[Textbook expanded form, Lean]
\label{thm:higgs-expanded}
\[
\mathcal{L}_{\rm Higgs}^{\rm pot}(\Phi)
= \lambda_{\rm eff}\,|\Phi|^4
- 2\lambda_{\rm eff}\,\texttt{lockinVev}^2\,|\Phi|^2
+ \lambda_{\rm eff}\,\texttt{lockinVev}^4 .
\]
This is the \hyperref[lean:higgs-expanded]{expanded Higgs potential identity}
(proved by \texttt{unfold} plus \texttt{ring}).
\end{theorem}

\begin{remark}[Lock-in calibration of the vev]
\label{rem:lockin}
The calibration $\texttt{lockinVev}^2=\eta_{\rm paper}$ when
$\Omega_k(m_{\rm lockin};m_{\rm lockin})=1$ is the
\hyperref[lean:lockin-vev-sq]{lock-in vev squared identity}; the curvature
ratio at lock-in is the
\hyperref[lean:omega-k-lockin]{lock-in $\Omega_k$ calibration}. The lock-in
vev is therefore \emph{not} a fitted constant: it is the square root of the
proved baryogenesis $\eta$-anchor, packaged once and reused unchanged in the
Higgs sector.
\end{remark}

\paragraph{Reading.}
The textbook Higgs page becomes, in HQIV notation, the diagonal scalar-norm
sum (kinetic) plus the expanded Higgs potential with
$(\lambda,v^2)=(\lambda_{\rm eff},\eta_{\rm paper})$. The constant
$\lambda_{\rm eff}\,v^4$ is a global vacuum offset that does not enter
equations of motion.

\section{Yukawa sector from the resonance ladder}
\label{sec:yukawa}

The Yukawa block of $\mathcal{L}_{\rm SM}$ is
$\mathcal{L}_Y=-\sum_f y_f\,\bar f\Phi f + \text{h.c.}$. In HQIV the Yukawa
coupling for each SM flavour is the textbook combination $y_f=\sqrt{2}\,
m_f/v$ with $v=\texttt{lockinVev}$ and $m_f$ \emph{not} a free parameter:

\[
m_f = m_f^{\mathrm{geom}}
= m_\tau^{\rm Pl}\;\big/\;\mathrm{resonanceProduct}\bigl(\mathrm{genIndex}(f)\bigr) ,
\]
where $m_f^{\mathrm{geom}}$ is the
\hyperref[lean:sm-mass-label]{geometric SM mass functional} and
$\mathrm{resonanceProduct}$ is built from the two
\hyperref[lean:resonance-steps]{resonance steps} of the charged-lepton
ladder: generations $\{8_v, 8_{s^+}, 8_{s^-}\}$ get product
$k_{\tau\mu}\,k_{\mu e}$, $k_{\tau\mu}$, $1$ respectively.

\begin{definition}[SM Yukawa coupling]
\label{def:yukawa}
For an SM flavour label $\ell\in\texttt{SMMassLabel}$ (the
\hyperref[lean:sm-mass-label-type]{SM mass-label type}),
\[
y_{\rm SM}(\ell)
:= \frac{\sqrt{2}\;m_\ell^{\mathrm{geom}}}{\texttt{lockinVev}}
\]
(the \hyperref[lean:y-SM]{SM Yukawa coupling}). For a symbolic Dirac bilinear
$B:\texttt{SMMassLabel}\to\mathbb{R}$,
\[
\mathcal{L}^{\rm SM}_{\rm Yukawa}(\ell, B)
:= -\,y_{\rm SM}(\ell)\;B(\ell)
\]
(the \hyperref[lean:L-SM-Yukawa]{per-flavour Yukawa term}); the all-flavour
aggregate is the list sum over twelve SM flavours (the
\hyperref[lean:all-flavours]{all-flavours list}, length certified by the
\hyperref[lean:all-flavours-len]{all-flavours length lemma}).
\end{definition}

\begin{theorem}[Yukawa identity, Lean]
\label{thm:yukawa-identity}
For every SM flavour $\ell$ and every $v=\texttt{lockinVev}\neq 0$,
\[
\frac{y_{\rm SM}(\ell)\;v}{\sqrt{2}} = m_\ell^{\mathrm{geom}}
= m_\tau^{\rm Pl}\;/\;\mathrm{resonanceProduct}(\mathrm{genIndex}(\ell)).
\]
This is the pair comprising the
\hyperref[lean:yukawa-mass]{Yukawa--mass identity} and the
\hyperref[lean:yukawa-resonance]{Yukawa resonance witness}.
\end{theorem}

\begin{theorem}[Yukawa linearity over flavours, Lean]
\label{thm:yukawa-sum}
For any flavour list $L$ and bilinear assignment $B$,
\[
\sum_{\ell\in L} \mathcal{L}^{\rm SM}_{\rm Yukawa}(\ell, B(\ell))
= -\,\sum_{\ell\in L} y_{\rm SM}(\ell)\,B(\ell) .
\]
This is the \hyperref[lean:yukawa-sum]{Yukawa sum identity}.
\end{theorem}

\begin{remark}[No fitted Yukawa matrix]
\label{rem:no-fits}
The textbook SM stores $\mathcal{O}(20)$ real Yukawa and mixing parameters
fitted from experiment. In HQIV these are replaced by the single rule
$y_f=\sqrt{2}\,m_f/v$ with $m_f$ from the discrete resonance ladder and
$\tau$-Planck anchor---not a simultaneous PDG mass solve. Laboratory masses
enter only as comparison targets under the proton lock-in witness discipline
of \cite{hqiv-mass-spectrum-paper,hqiv-kirchhoff-paper}. Mixing structure is
the triality-orbit branching of the SM embedding module; hierarchies are the
harmonic ladder, not a tuned Yukawa table.
\end{remark}

\paragraph{Reading.}
The textbook Yukawa ``page'' $-y_e\bar L\Phi e_R - y_u\bar Q\tilde\Phi u_R
- y_d\bar Q\Phi d_R + \text{h.c.}$ (three generations) becomes, in HQIV
notation, a list-sum over twelve flavours each weighted by
$\sqrt{2}\,m_\tau^{\rm Pl}/(\texttt{lockinVev}\cdot\mathrm{resonanceProduct})$.
\emph{Three} numbers ($m_\tau^{\rm Pl}$, $k_{\tau\mu}$, $k_{\mu e}$) fix the
entire Yukawa sector.

\section{Assembled SM Lagrangian and headline theorems}
\label{sec:assembly}

\begin{definition}[SM Lagrangian record (Lean)]
\label{def:SM-record}
The \hyperref[lean:SM-record]{SM Lagrangian record} carries five real-valued fields
$(\mathcal{L}_{\rm gauge},\mathcal{L}_{\rm fermion},
\mathcal{L}_{\rm Higgs}^{\rm kin},\mathcal{L}_{\rm Higgs}^{\rm pot},
\mathcal{L}_Y)$. The total is
\[
\mathcal{L}_{\rm SM}^{\rm total}
:= \mathcal{L}_{\rm gauge} + \mathcal{L}_{\rm fermion}
+ \mathcal{L}_{\rm Higgs}^{\rm kin} - \mathcal{L}_{\rm Higgs}^{\rm pot}
+ \mathcal{L}_Y
\]
(the \hyperref[lean:SM-total]{SM Lagrangian total}).
\end{definition}

\begin{definition}[SM Lagrangian from HQIV (Lean)]
\label{def:SM-fromHQIV}
Given an HQIV gauge potential $A$, a triality-indexed current family $J$,
an octonion scalar $\Phi$, a per-direction scalar slot $\Phi_{\rm dir}$, and
a flavour bilinear assignment $B$, the
\hyperref[lean:SM-fromHQIV]{from-HQIV constructor} sets
\[
\mathcal{L}_{\rm gauge} = L_O^{\rm kin}(A),\quad
\mathcal{L}_{\rm fermion} = \sum_{\rm gen}L_O^{\rm src}(J(\rm gen),A),
\]
\[
\mathcal{L}_{\rm Higgs}^{\rm kin} = \sum_\mu \mathrm{scalarNormSq}(\Phi_{\rm dir}(\mu)),\quad
\mathcal{L}_{\rm Higgs}^{\rm pot} = \texttt{higgsPotential}(\lambda_{\rm eff},\texttt{lockinVev},\Phi),
\]
\[
\mathcal{L}_Y = \!\!\sum_{\ell\in\texttt{all\_SM\_flavours}}\!\! \mathcal{L}^{\rm SM}_{\rm Yukawa}(\ell, B(\ell)) .
\]
\end{definition}

\begin{theorem}[Headline: SM Lagrangian = HQIV discrete action]
\label{thm:headline}
For every $(A, J, \Phi, \Phi_{\rm dir}, B)$, writing $L$ for the from-HQIV
record at those arguments,
\begin{enumerate}[leftmargin=2em]
\item $\mathcal{L}_{\rm gauge}(L) = L_O^{\rm kin}(A)$;
\item $\mathcal{L}_{\rm fermion}(L) = L_O^{\rm src}\!\left(J_{\rm SM}^{\rm tot}(J),A\right)$,
where $J_{\rm SM}^{\rm tot}$ is the \hyperref[lean:J-SM-total]{triality-summed SM current};
\item $\mathcal{L}_{\rm Higgs}^{\rm kin}(L)
= \sum_\mu \mathrm{scalarNormSq}(\Phi_{\rm dir}(\mu))$;
\item $\mathcal{L}_{\rm Higgs}^{\rm pot}(L)
= \texttt{higgsPotential}(\lambda_{\rm eff},\texttt{lockinVev},\Phi)$;
\item $\mathcal{L}_Y(L)$ is the twelve-flavour Yukawa sum with weights
$\sqrt{2}\,m_\ell^{\mathrm{geom}}/\texttt{lockinVev}$.
\end{enumerate}
This is the \hyperref[lean:headline]{SM-from-HQIV discrete-action theorem}.
\end{theorem}

\begin{theorem}[Total HQIV action split, Lean]
\label{thm:total-split-lean}
With $J_{\rm SM}^{\rm tot}(J)$ as in Theorem~\ref{thm:headline},
\begin{multline*}
\mathcal{S}_{\mathrm{HQIV}}^{\mathrm{total}}\bigl(J_{\rm SM}^{\rm tot}(J),A,\phi_{\rm val},\rho_m,\rho_r\bigr)
\\
= \biggl(L_O^{\rm kin}(A)
+ 4\pi\sum_{\rm gen} L_O^{\rm src}(J({\rm gen}),A)
+ L_O^\phi(A,\phi_{\rm val})\biggr) + S_{\rm HQVM\,grav}(\phi_{\rm val},\rho_m,\rho_r).
\end{multline*}
This is the \hyperref[lean:total-split]{total-action split theorem}; its
proof uses two invocations of the source-coupling additivity lemma, then
\texttt{ring} on the residue.
\end{theorem}

\paragraph{Reading.}
Theorem~\ref{thm:total-split-lean} exposes the dictionary
\[
\boxed{\;
\mathcal{S}_{\rm HQIV}
= \mathcal{L}_{\rm gauge}^{\rm SM}
+ \frac{1}{4\pi}\mathcal{L}_{\rm fermion}^{\rm SM}
+ \mathcal{L}_{\rm Higgs\;cross}^{\rm SM}
+ S_{\rm HQVM\,grav}
\;}
\]
between the page-long SM Lagrangian and the HQIV total cell. The
$1/(4\pi)$ factor is the discrete normalization from the Maxwell-source
convention of the \hyperref[lean:LO-maxwell]{general O-Maxwell Lagrangian};
the Higgs-cross label captures the $\phi$-coupling remnant once $\Phi$ is
collapsed to the lattice scalar at the electroweak shell (shell index $5$ in
the SM--GR module).

\section{Parameter accounting and witnesses}
\label{sec:parameters}

The full from-HQIV construction consumes the following \emph{derived}
discrete inputs. None is a simultaneous PDG fit; laboratory comparison uses
the proton lock-in witness discipline of
\cite{hqiv-mass-spectrum-paper,hqiv-kirchhoff-paper}.

\begin{table}[ht]
\caption{Discrete parameter inputs to the HQIV-derived SM Lagrangian. The
$\tau$-Planck entry is the lepton-scale anchor on the resonance ladder (not a
PDG pole mass imported into the proof).}
\label{tab:SM-parameters}
\centering
\footnotesize
\begin{tabular}{@{}p{4.0cm}p{5.0cm}p{4.4cm}@{}}
\toprule
Constant & HQIV source & Appendix link \\
\midrule
$\alpha=3/5$ & Lattice stars-and-bars normalization &
  \hyperref[lean:alpha-35]{lattice $\alpha$ identity} \\
$\gamma_{\rm HQIV}=2/5$ & Horizon split $\gamma=1-\alpha$ &
  \hyperref[lean:gamma-25]{horizon-split $\gamma$ identity} \\
$\alpha_{\rm GUT}=1/42$ & Cube axes $\times$ octonion imaginaries &
  \hyperref[lean:alpha-gut]{GUT coupling identity} \\
$\lambda_{\rm eff}$ & Effective Higgs quartic at EW shell &
  \hyperref[lean:lambda-eff]{effective quartic} \\
$v$ (lock-in vev) & $\sqrt{\eta_{\rm paper}\,\Omega_k(M;M)}$ &
  \hyperref[lean:lockin-vev]{lock-in vev} \\
$m_\tau^{\rm Pl}$ & $\tau$ scale on discrete ladder (not PDG) &
  \hyperref[lean:m-tau-pl]{tau Planck mass} \\
$k_{\tau\mu}$, $k_{\mu e}$ & Charged-lepton resonance steps &
  \hyperref[lean:resonance-steps]{two resonance steps} \\
\bottomrule
\end{tabular}
\end{table}

\begin{theorem}[Parameter witness, Lean]
\label{thm:param-witness}
$\alpha=3/5$, $\gamma_{\rm HQIV}=2/5$, $\alpha_{\rm GUT}=1/42$
simultaneously. This is the \hyperref[lean:param-witness]{parameter witness},
discharged by the three rational identities listed in the table.
\end{theorem}

\begin{theorem}[Yukawa resonance witness, Lean]
\label{thm:yukawa-resonance-witness}
For every SM flavour $\ell$, when $\texttt{lockinVev}\neq 0$,
\[
\frac{y_{\rm SM}(\ell)\,\texttt{lockinVev}}{\sqrt{2}}
= m_\tau^{\rm Pl}\;/\;\mathrm{resonanceProduct}(\mathrm{genIndex}(\ell)).
\]
This is the \hyperref[lean:yukawa-resonance]{Yukawa resonance witness}.
\end{theorem}

\begin{remark}[Comparison with textbook parameter count]
\label{rem:param-count}
The textbook SM has $\mathcal{O}(20)$ fitted real parameters. The HQIV
bridge compresses them to: three lattice rationals, one $\tau$-Planck ladder
anchor, two resonance steps, and one effective quartic---all discrete outputs
of the upstream spine. Gauge couplings trace to $\alpha_{\rm GUT}$ plus the
$\phi$-correction of \cite{hqiv-oct-action-paper}; charged-fermion masses
trace to the resonance ladder; CKM/PMNS structure traces to triality
branching; $(\mu^2,\lambda)$ trace to $(\eta_{\rm paper},\lambda_{\rm eff})$.
PDG values are comparison targets under the single-witness policy, not rows
in the solve.
\end{remark}

\section{Outlook: what comes next}
\label{sec:outlook}

The structural bridge of Theorem~\ref{thm:headline} is exact at the
bookkeeping level. To translate it to textbook smooth-manifold notation (a
calculation-approximation layer for comparison with
\cite{peskin1995,schwartz2014}, \emph{not} a foundational completion of HQIV),
three Tier-III tasks remain (ordering consistent with
\cite{hqiv-oct-action-paper}):

\begin{enumerate}[leftmargin=2em]
\item \textbf{Variational fermion sector.} Promote the symbolic current to a
  typed Dirac bilinear on the Cl$(0,6)$ spinor carrier already present in the
  library; the abelian source coupling of Section~\ref{sec:fermion} becomes
  the textbook minimal coupling.
\item \textbf{Non-abelian gauge kinetic.} Replace abelian $F^a{}_{\mu\nu}$ in
  the weak and strong sectors by the non-abelian weak field-strength scaffold
  and the $\mathrm{SU}(3)_c$ chart Lie law; the reduction theorem upstream
  guarantees collapse to the abelian core when the commutator vanishes.
\item \textbf{Higgs covariant derivative.} Wire full gauge connection data into
  the fermion and Higgs slots; verify the kinetic shadow reduces to the
  scalar-norm sum in the abelian limit.
\end{enumerate}

None of these tasks requires new physics input. The discrete constants, vev,
and mass ladder are already proved; smooth-manifold form is a translation
layer only---not a continuum-completion programme.

\section*{Acknowledgments}
AI-assisted drafting and repository tooling were used during manuscript
preparation. Formal claims refer to the HQIV-LEAN library \cite{hqiv-lean};
the headline module builds clean under \texttt{lake build}
\texttt{Hqiv.Physics.StandardModelLagrangianFromDiscreteAction}.

\appendix
\section{Lean module index}
\label{app:lean-index}

The Lean development lives at \url{https://github.com/HQIV/hqiv-lean}
\cite{hqiv-lean}. Reader-friendly anchors in the body link here.

\begingroup
\sloppy
\setlength{\emergencystretch}{4em}
\begin{description}[leftmargin=0em,style=nextline]
\item[\phantomsection\label{lean:sm-lagrangian-module}\textbf{Standard-Model Lagrangian bridge.}]
  Headline module: five sector equalities and parameter witness.
  \leanfile{Hqiv/Physics/StandardModelLagrangianFromDiscreteAction.lean}.
\item[\phantomsection\label{lean:action}\textbf{O-Maxwell action module.}]
  Discrete gauge action upstream. \leanfile{Hqiv/Physics/Action.lean}.
\item[\phantomsection\label{lean:F-from-A}\textbf{Field-strength constructor.}]
  $F^a{}_{\mu\nu}(A)=A^a{}_\nu-A^a{}_\mu$. \leanid{F\_from\_A}.
\item[\phantomsection\label{lean:LO-kin}\textbf{Octonion kinetic density.}]
  $L_O^{\rm kin}$. \leanid{L\_O\_kinetic}.
\item[\phantomsection\label{lean:LO-src}\textbf{Octonion source coupling.}]
  $L_O^{\rm src}$. \leanid{L\_O\_source\_general}.
\item[\phantomsection\label{lean:LO-phi}\textbf{Scalar $\phi$-coupling slot.}]
  \leanid{L\_O\_phi\_coupling}.
\item[\phantomsection\label{lean:LO-maxwell}\textbf{General O-Maxwell Lagrangian.}]
  \leanid{L\_O\_Maxwell\_general}.
\item[\phantomsection\label{lean:SHQVM-grav}\textbf{HQVM gravitational constraint.}]
  \leanid{S\_HQVM\_grav}.
\item[\phantomsection\label{lean:friedmann-iff}\textbf{Friedmann equivalence.}]
  \leanid{S\_HQVM\_grav\_zero\_iff\_Friedmann}.
\item[\phantomsection\label{lean:action-total}\textbf{HQIV total action.}]
  \leanid{action\_total\_general}.
\item[\phantomsection\label{lean:forces}\textbf{Forces module.}]
  \leanfile{Hqiv/Physics/Forces.lean}.
\item[\phantomsection\label{lean:O-to-sector}\textbf{Force-sector map.}]
  \leanid{O\_component\_to\_sector}.
\item[\phantomsection\label{lean:O-zero-EM}\textbf{Zeroth component is EM.}]
  \leanid{O\_component\_zero\_is\_EM}.
\item[\phantomsection\label{lean:LO-kin-sector}\textbf{Sector-projected kinetic.}]
  \leanid{L\_O\_kinetic\_sector}.
\item[\phantomsection\label{lean:LO-kinetic-split}\textbf{Three-sector kinetic split.}]
  \leanid{L\_O\_kinetic\_eq\_sum\_of\_sector\_pieces}.
\item[\phantomsection\label{lean:LSM-B-kin}\textbf{Hypercharge kinetic aggregate.}]
  \leanid{L\_SM\_B\_kinetic}.
\item[\phantomsection\label{lean:LSM-W-kin}\textbf{Weak kinetic aggregate.}]
  \leanid{L\_SM\_W\_kinetic}.
\item[\phantomsection\label{lean:LSM-G-kin}\textbf{Strong kinetic aggregate.}]
  \leanid{L\_SM\_G\_kinetic}.
\item[\phantomsection\label{lean:gauge-eq-kin}\textbf{Gauge--kinetic equality.}]
  \leanid{L\_SM\_gauge\_kinetic\_eq\_L\_O\_kinetic}.
\item[\phantomsection\label{lean:sm-embedding}\textbf{SM embedding module.}]
  \leanfile{Hqiv/Algebra/SMEmbedding.lean}.
\item[\phantomsection\label{lean:fermion-minimal}\textbf{Per-generation minimal coupling.}]
  \leanid{L\_SM\_fermion\_minimal\_coupling}.
\item[\phantomsection\label{lean:three-gen-reps}\textbf{Three generations from triality.}]
  \leanid{three\_generations\_from\_triality\_reps}.
\item[\phantomsection\label{lean:three-gen}\textbf{Three-generation equality.}]
  \leanid{L\_SM\_three\_generations\_eq\_total\_source}.
\item[\phantomsection\label{lean:J-add}\textbf{Source-coupling additivity.}]
  \leanid{L\_O\_source\_general\_add\_J}.
\item[\phantomsection\label{lean:sm-gr}\textbf{SM--GR unification module.}]
  \leanfile{Hqiv/Physics/SM\_GR\_Unification.lean}.
\item[\phantomsection\label{lean:sm-gen-index}\textbf{Generation index map.}]
  \leanid{smGenerationIndex}.
\item[\phantomsection\label{lean:lepton-resonance}\textbf{Charged-lepton resonance.}]
  \leanfile{Hqiv/Physics/ChargedLeptonResonance.lean}.
\item[\phantomsection\label{lean:resonance-product}\textbf{Resonance-product factor.}]
  \leanid{resonanceProduct}.
\item[\phantomsection\label{lean:weak-higgs}\textbf{Weak/Higgs scaffold.}]
  \leanfile{Hqiv/Physics/WeakHiggsFromOMaxwellScaffold.lean}.
\item[\phantomsection\label{lean:scalar-norm}\textbf{Octonion scalar norm.}]
  \leanid{scalarNormSq}.
\item[\phantomsection\label{lean:higgs-pot}\textbf{Higgs potential.}]
  \leanid{higgsPotential}.
\item[\phantomsection\label{lean:lockin-vev}\textbf{Lock-in vev.}]
  \leanid{lockinVev}.
\item[\phantomsection\label{lean:lambda-eff}\textbf{Effective quartic.}]
  \leanid{lambda\_eff}.
\item[\phantomsection\label{lean:LSM-Higgs-kin}\textbf{SM Higgs kinetic shadow.}]
  \leanid{L\_SM\_Higgs\_kinetic}.
\item[\phantomsection\label{lean:LSM-Higgs-pot}\textbf{SM Higgs potential shadow.}]
  \leanid{L\_SM\_Higgs\_potential}.
\item[\phantomsection\label{lean:higgs-expanded}\textbf{Expanded Higgs potential.}]
  \leanid{L\_SM\_Higgs\_potential\_expanded}.
\item[\phantomsection\label{lean:lockin-vev-sq}\textbf{Lock-in vev squared.}]
  \leanid{lockinVev\_sq\_eq\_eta\_paper}.
\item[\phantomsection\label{lean:omega-k-lockin}\textbf{Lock-in $\Omega_k$ calibration.}]
  \leanid{omega\_k\_lockin\_calibration}.
\item[\phantomsection\label{lean:sm-mass-label}\textbf{Geometric SM mass functional.}]
  \leanid{smMassFromGeometryLabel}.
\item[\phantomsection\label{lean:sm-mass-label-type}\textbf{SM mass-label type.}]
  \leanid{SMMassLabel}.
\item[\phantomsection\label{lean:resonance-steps}\textbf{Two resonance steps.}]
  \leanid{resonance\_k\_tau\_mu}, \leanid{resonance\_k\_mu\_e}.
\item[\phantomsection\label{lean:m-tau-pl}\textbf{Tau Planck mass.}]
  \leanid{m\_tau\_Pl}.
\item[\phantomsection\label{lean:y-SM}\textbf{SM Yukawa coupling.}]
  \leanid{y\_SM}.
\item[\phantomsection\label{lean:L-SM-Yukawa}\textbf{Per-flavour Yukawa term.}]
  \leanid{L\_SM\_Yukawa\_flavour}.
\item[\phantomsection\label{lean:all-flavours}\textbf{All-flavours list.}]
  \leanid{all\_SM\_flavours}.
\item[\phantomsection\label{lean:all-flavours-len}\textbf{All-flavours length.}]
  \leanid{all\_SM\_flavours\_length}.
\item[\phantomsection\label{lean:yukawa-mass}\textbf{Yukawa--mass identity.}]
  \leanid{y\_SM\_times\_v\_over\_sqrt2\_eq\_mass}.
\item[\phantomsection\label{lean:yukawa-resonance}\textbf{Yukawa resonance witness.}]
  \leanid{SM\_Lagrangian\_yukawa\_resonance\_witness}.
\item[\phantomsection\label{lean:yukawa-sum}\textbf{Yukawa sum identity.}]
  \leanid{L\_SM\_Yukawa\_sum\_eq\_sum}.
\item[\phantomsection\label{lean:SM-record}\textbf{SM Lagrangian record.}]
  \leanid{SM\_Lagrangian}.
\item[\phantomsection\label{lean:SM-fromHQIV}\textbf{From-HQIV constructor.}]
  \leanid{SM\_Lagrangian.fromHQIV}.
\item[\phantomsection\label{lean:SM-total}\textbf{SM Lagrangian total.}]
  \leanid{SM\_Lagrangian.total}.
\item[\phantomsection\label{lean:headline}\textbf{SM-from-HQIV discrete-action theorem.}]
  \leanid{SM\_Lagrangian\_from\_HQIV\_discrete\_action}.
\item[\phantomsection\label{lean:J-SM-total}\textbf{Triality-summed SM current.}]
  \leanid{J\_SM\_total}.
\item[\phantomsection\label{lean:total-split}\textbf{Total-action split theorem.}]
  \leanid{HQIV\_total\_action\_eq\_SM\_gauge\_fermion\_plus\_phi\_coupling\_plus\_grav}.
\item[\phantomsection\label{lean:param-witness}\textbf{Parameter witness.}]
  \leanid{SM\_Lagrangian\_parameter\_witness}.
\item[\phantomsection\label{lean:alpha-35}\textbf{Lattice $\alpha$ identity.}]
  \leanid{alpha\_eq\_3\_5}.
\item[\phantomsection\label{lean:gamma-25}\textbf{Horizon-split $\gamma$ identity.}]
  \leanid{gamma\_eq\_2\_5}.
\item[\phantomsection\label{lean:alpha-gut}\textbf{GUT coupling identity.}]
  \leanid{alpha\_GUT\_eq\_1\_42}.
\item[\phantomsection\label{lean:proton-lockin}\textbf{Proton lock-in scale-witness.}]
  \leanid{ScaleWitness.proton\_lockin} in
  \leanfile{Hqiv/Physics/ScaleWitness.lean}.
\end{description}
\endgroup

\bibliographystyle{plain}
\bibliography{../references}

\end{document}
